\documentclass[11pt]{article}
\usepackage[margin=1in]{geometry}
\usepackage[T1]{fontenc}
\usepackage[utf8]{inputenc}
\usepackage{lmodern}
\usepackage{microtype}
\usepackage{xcolor}
\usepackage{amsmath,amssymb,amsthm,mathtools}
\usepackage{bm}
\usepackage{hyperref}
\usepackage{enumitem}
\usepackage{booktabs}
\usepackage{array}
\usepackage{listings}
\usepackage{inconsolata}
\usepackage{fancyhdr}
\usepackage{titlesec}

\hypersetup{
  colorlinks=true,
  linkcolor=blue!50!black,
  urlcolor=blue!50!black,
  citecolor=blue!50!black,
  pdftitle={Middle Product via Multipairing},
  pdfauthor={OpenAI}
}

\pagestyle{fancy}
\fancyhf{}
\lhead{Middle Product via Multipairing}
\rhead{\thepage}
\setlength{\headheight}{14pt}

\titleformat{\section}{\large\bfseries\color{blue!35!black}}{\thesection.}{0.5em}{}
\titleformat{\subsection}{\normalsize\bfseries\color{blue!20!black}}{\thesubsection.}{0.5em}{}

\lstdefinestyle{mystyle}{
  basicstyle=\ttfamily\small,
  keywordstyle=\bfseries\color{blue!50!black},
  commentstyle=\itshape\color{black!55},
  frame=single,
  framerule=0.4pt,
  rulecolor=\color{black!20},
  backgroundcolor=\color{black!2},
  showstringspaces=false,
  columns=fullflexible,
  keepspaces=true,
  xleftmargin=0.5em,
  xrightmargin=0.5em
}
\lstset{style=mystyle}

\newcommand{\GT}{G_T}

\begin{document}

\begin{center}
  {\LARGE \textbf{Middle Product via Multipairing}}\\[0.5em]
  {\large Technical note}\\[0.75em]
\end{center}

\noindent
This note describes a direct middle-product implementation based on
multipairing. The goal is not to replace the optimized FFT formulation at
large sizes. The goal is to understand a small-size execution path whose
native batching primitive is the pairing itself.

\section{Problem statement}

Let
\[
A(x) = a_0 + a_1 x + \cdots + a_{n-1} x^{n-1}
\]
and
\[
B(x) = b_0 + b_1 x + \cdots + b_{2n-1} x^{2n-1},
\]
where the coefficients of $A$ lie in one pairing input group, the coefficients
of $B$ lie in the other pairing input group, and coefficient multiplication is
the bilinear pairing
\[
a_i \circ b_j \in \GT.
\]

Define
\[
C(x) = A(x)B(x) = \sum_{k=0}^{3n-2} C_k x^k.
\]
We only want the middle window
\[
C_n, C_{n+1}, \dots, C_{2n-1}.
\]

For each $r \in \{0,\dots,n-1\}$,
\[
C_{n+r} = \sum_{i=0}^{n-1} a_i \circ b_{n+r-i}.
\]
This is the middle product.

\section{Naive direct evaluation}

The literal implementation is:

\begin{lstlisting}
for r in 0..n-1:
    acc = 0 in G_T
    for i in 0..n-1:
        acc += pair(a[i], b[n + r - i])
    out[r] = acc
\end{lstlisting}

This performs $n^2$ bilinear products.
It is simple and exact, but it does not exploit the fact that pairings are
normally implemented as a two-phase computation:
\begin{enumerate}[itemsep=0.3em, topsep=0.4em]
  \item Miller accumulation,
  \item final exponentiation.
\end{enumerate}

That structure motivates a direct multipairing formulation.

\section{Multipairing reformulation}

For a fixed output index $r$, define the pair list
\[
(a_0, b_{n+r}),
\ (a_1, b_{n+r-1}),
\ \dots,
\ (a_{n-1}, b_r).
\]
Then the same coefficient can be written as one multipairing:
\[
C_{n+r}
=
\mathrm{MultiPair}\big((a_i, b_{n+r-i})_{i=0}^{n-1}\big).
\]

This does \emph{not} change the asymptotic cost.
We still have $n$ outputs, each involving $n$ input pairs, so total work
remains $O(n^2)$.
The gain is in constant factors:
\begin{itemize}[itemsep=0.3em, topsep=0.4em]
  \item one Miller pass can accumulate many pairs,
  \item one final exponentiation serves the whole output,
  \item line-table precomputation can be reused,
  \item zero and trivial terms can be skipped in bulk.
\end{itemize}

\section{Why multipairing is the natural primitive}

MSM is the right batching primitive for linear combinations inside $G_1$ or
$G_2$.
This problem is different.
The coefficient formula is already bilinear in $(a_i, b_j)$ and lands in
\GT.

So for the direct middle-product path, the natural batching primitive is:
\begin{quote}
Batch many pair contributions into one multipairing per output coefficient.
\end{quote}

This keeps the computation aligned with the algebra:
\[
\sum_i a_i \circ b_{n+r-i}.
\]
It also matches what pairing libraries are usually optimized to do well:
share Miller-loop work across many pairs and postpone the final exponentiation
until the end of the batch.

\section{Precomputation strategy}

The most useful precomputation is on the $B$ side.
For each coefficient $b_j$, precompute the data needed to evaluate a Miller
contribution quickly.
In practice this means storing one of:
\begin{itemize}[itemsep=0.3em, topsep=0.4em]
  \item an affine or other reduced representation of $b_j$,
  \item a prepared line table for $b_j$,
  \item any backend-specific pairing precompute object.
\end{itemize}

Then output $r$ only needs the contiguous prepared window
\[
b_r,\ b_{r+1},\ \dots,\ b_{r+n-1},
\]
read in reverse against
\[
a_{n-1},\ a_{n-2},\ \dots,\ a_0.
\]

A useful precompute object therefore stores:
\begin{enumerate}[itemsep=0.3em, topsep=0.4em]
  \item the original problem size $n$,
  \item prepared data for all $2n$ coefficients of $B$,
  \item optional zero flags for fast skipping,
  \item optional contiguous views or index helpers for output windows.
\end{enumerate}

\section{Concrete direct algorithm}

\subsection{Offline precomputation}

\begin{lstlisting}
precompute_B_for_multipair(B):
    pre.n = len(B) / 2
    for j in 0..2*n-1:
        pre.zero[j] = is_zero(B[j])
        if not pre.zero[j]:
            pre.lines[j] = prepare_pairing_lines(B[j])
    return pre
\end{lstlisting}

\subsection{Online evaluation}

\begin{lstlisting}
middle_product_multipair(A, pre):
    n = len(A)
    for r in 0..n-1:
        mp = multipair_identity()
        for i in 0..n-1:
            j = n + r - i
            if is_zero(A[i]) or pre.zero[j]:
                continue
            multipair_accumulate(mp, A[i], pre.lines[j])
        out[r] = multipair_finalize(mp)
    return out
\end{lstlisting}

This is still the same exact formula
\[
C_{n+r} = \sum_{i=0}^{n-1} a_i \circ b_{n+r-i},
\]
just implemented through the library's batched pairing machinery rather than
as $n$ separate pairings.

\section{Update and combine view}

If the API already separates the middle product into:
\begin{enumerate}[itemsep=0.3em, topsep=0.4em]
  \item a main middle-product evaluation,
  \item a constant-term or shifted update,
  \item a combine step in \GT,
\end{enumerate}
then the direct multipairing path fits that shape naturally.

For example, the contribution of $a_0$ alone is
\[
a_0 \circ b_n,\ a_0 \circ b_{n+1},\ \dots,\ a_0 \circ b_{2n-1},
\]
which is just the length-$n$ vector obtained by pairing one fixed $a_0$ with
the upper half of $B$.
The remaining coefficients can be treated similarly, and the final result is
the accumulation of those per-coefficient contributions in \GT.

This observation does not change complexity, but it means the direct
multipairing path can often be made to mirror the same high-level API shape as
the FFT path.

\section{Comparison with the optimized FFT path}

The optimized FFT middle-product path has much better asymptotic behavior:
\[
O(n \log n)
\]
group transforms, plus pointwise pairing work in the evaluation domain.

The direct multipairing path remains:
\[
O(n^2)
\]
in total pair inputs.

So the expected tradeoff is:
\begin{itemize}[itemsep=0.3em, topsep=0.4em]
  \item \textbf{Multipairing path:} small sizes, simple control flow, good
  constant factors, minimal transform machinery.
  \item \textbf{FFT path:} medium and large sizes, better asymptotics, more
  structure and precomputation.
\end{itemize}

The direct path is therefore best viewed as a small-size competitor or fallback,
not as a replacement for the FFT design.

\section{Benchmark plan}

The first benchmark goal is very simple:
\begin{quote}
Find the largest $n$ for which direct multipairing beats the optimized FFT
middle product on real AVX hardware.
\end{quote}

That benchmark should measure at least:
\begin{enumerate}[itemsep=0.3em, topsep=0.4em]
  \item direct multipairing with no precompute,
  \item direct multipairing with prepared $B$ data,
  \item optimized FFT with its existing precompute path,
  \item online-only versus online-plus-offline totals.
\end{enumerate}

Representative sizes might be
\[
n \in \{2,4,8,16,32,64,128\}.
\]

The outcome of this benchmark is a dispatch threshold:
\begin{itemize}[itemsep=0.3em, topsep=0.4em]
  \item use direct multipairing below the threshold,
  \item use the optimized FFT path above the threshold.
\end{itemize}

\section{Implementation plan}

\subsection{Phase 1: standalone batched pairing kernel}

Build a reusable kernel for:
\begin{enumerate}[itemsep=0.3em, topsep=0.4em]
  \item preparing one $B$ coefficient for pairing,
  \item accumulating many $(A_i, B_j)$ contributions into one multipair object,
  \item finalizing the batched result into \GT.
\end{enumerate}

This should be benchmarked independently before wiring it into a full
middle-product path.

\subsection{Phase 2: direct middle-product implementation}

Build a clean direct implementation around that kernel:
\begin{itemize}[itemsep=0.3em, topsep=0.4em]
  \item precompute $B$,
  \item for each output $r$, accumulate the diagonal
  $(a_i, b_{n+r-i})$ into one multipair,
  \item finalize to produce $C_{n+r}$.
\end{itemize}

\subsection{Phase 3: correctness}

Compare the direct multipairing outputs against:
\begin{enumerate}[itemsep=0.3em, topsep=0.4em]
  \item a naive reference implementation,
  \item the existing optimized FFT middle-product path.
\end{enumerate}

Use deterministic vectors and cover zero-heavy cases, boundary sizes, and
prepared-line versus unprepared modes.

\subsection{Phase 4: crossover measurement}

Benchmark the direct multipairing path against the optimized FFT path on real
AVX hardware.
The deliverable is an empirical threshold, not a proof.

\subsection{Phase 5: optional dispatch integration}

Only after the crossover is understood should any runtime dispatch be added.
The simplest policy is:
\begin{lstlisting}
if n <= threshold:
    use direct multipairing
else:
    use optimized FFT
\end{lstlisting}

\section{What to keep from the current work}

The current experimental work should be filtered rather than discarded
wholesale.
The important distinction is:
\begin{quote}
Keep the scaffolding and measurement structure; replace the naive per-term
pairing kernel.
\end{quote}

\subsection{Keep}

The following pieces remain useful for the multipairing plan:
\begin{itemize}[itemsep=0.3em, topsep=0.4em]
  \item the standalone planning note itself,
  \item the test wiring and benchmark wiring,
  \item deterministic input generation for $G_1$ and $G_2$ vectors,
  \item output-comparison helpers against the existing FFT path,
  \item the offline-versus-online timing split,
  \item the scalar-versus-AVX benchmark structure,
  \item the precompute / evaluate / update / combine API shape.
\end{itemize}

These pieces are useful because they are agnostic to whether the internal
kernel is naive direct pairing or true multipairing.

\subsection{Rewrite}

The current direct kernel should not be kept as the candidate implementation.
It should be replaced by a genuine multipairing core:
\begin{itemize}[itemsep=0.3em, topsep=0.4em]
  \item replace the one-term-at-a-time pairing helper,
  \item replace the nested ``pair then add into \GT'' inner loop,
  \item introduce a multipair object with accumulate and finalize steps,
  \item preserve prepared-line and zero-mask precomputation where useful.
\end{itemize}

In other words, the outer shape survives, but the computational heart changes.

\subsection{Demote}

The current naive direct implementation is still valuable as a reference
oracle, but it should no longer be treated as the performance candidate.
Any benchmark conclusions derived from:
\begin{quote}
naive direct versus FFT
\end{quote}
should be treated as obsolete for the multipairing decision.

\subsection{Practical repository guidance}

In repository terms, the intended reuse pattern is:
\begin{itemize}[itemsep=0.3em, topsep=0.4em]
  \item keep the tests and benchmarks as harnesses,
  \item keep the deterministic generators and comparison logic,
  \item keep the precompute container idea on the $B$ side,
  \item throw away the naive per-term pairing core,
  \item replace it with true multipair accumulation.
\end{itemize}

\section{Summary}

The direct middle-product formula is already pairing-shaped.
That makes multipairing the most natural small-size batching primitive.
It does not change the $O(n^2)$ asymptotics, but it can substantially reduce
constant factors by sharing Miller work and amortizing final exponentiation.

So the practical strategy is:
\begin{enumerate}[itemsep=0.3em, topsep=0.4em]
  \item keep the optimized FFT path as the large-size solution,
  \item build a direct multipairing path as a small-size competitor,
  \item benchmark both on real AVX hardware,
  \item dispatch based on measured crossover.
\end{enumerate}

\end{document}
